\documentclass{article}
\usepackage{neurips_2026}
% neurips_2026.sty loads natbib without options (defaults to author-year).
% Switch to numeric mode so our \bibitem list works without warnings.
\setcitestyle{numbers,sort&compress}

\usepackage{booktabs}
\usepackage{graphicx}
\usepackage{amsmath}
\usepackage{amssymb}
\usepackage{hyperref}

\title{scPTR: Decomposing Post-Transcriptional Regulation at Single-Cell Resolution}

\author{Anonymous Author(s) \\ Affiliation \\ Address \\ \texttt{email}}

\begin{document}

\maketitle

\begin{abstract}
Transcript abundance is jointly determined by mRNA synthesis and degradation, yet single-cell computational frameworks have addressed only the synthesis side. RNA velocity recovers transcriptional dynamics from spliced/unspliced ratios, but mRNA degradation---orchestrated by RNA-binding proteins (RBPs), microRNAs, and \textit{cis}-regulatory elements---remains an unobserved axis. We present \textbf{scPTR}, a framework that estimates per-cell, per-gene mRNA degradation rates ($\gamma$) from standard scRNA-seq, discovers cell subpopulations defined by post-transcriptional programs, and infers RBP regulatory networks at single-cell resolution. scPTR's $\gamma$ estimates correlate with published mRNA half-lives (Spearman $\rho{=}{-}0.81$ on sci-fate metabolic-labeling data), outperforming scVelo steady-state ($\rho{=}{-}0.37$) and velVI ($\rho{=}{-}0.28$); 59\% of 215 miRNA target families are significantly enriched ($p{=}4.7{\times}10^{-65}$). Clustering in $\gamma$-space reveals \textit{expression-invisible} post-transcriptional states in 5 of 8 pancreatic and 6 of 11 hippocampal cell types. Post-transcriptional velocity is orthogonal to RNA velocity and shows that degradation-rate changes \textit{precede} expression changes for 63\% of transition genes. Library-size-corrected RBP--target inference recovers known cancer dependencies. We additionally introduce \textsc{DeepPTR}, a structured VAE with a kinetic decoder that factorises the latent space into transcriptional and post-transcriptional components under a negative binomial observation model. Code, data, and notebooks reproducing every figure are released anonymously with the supplementary material.
\end{abstract}

\section{Introduction}

Steady-state mRNA abundance is set by the balance of synthesis and degradation. Single-cell computational biology has invested heavily in the synthesis side: RNA velocity \cite{lamanno2018} recovers transcriptional dynamics from the ratio of nascent (unspliced) to mature (spliced) reads, and modern deep-generative variants \cite{gayoso2024,gao2022} now infer latent dynamical states under expressive priors. Degradation has remained out of view, despite its established role in setting cell identity and fate \cite{hentze2018}. RNA-binding proteins (RBPs), microRNAs, and 3' UTR \textit{cis}-elements jointly determine transcript half-lives that vary by more than three orders of magnitude across the transcriptome and shift sharply during cell-state transitions, but no method to date estimates per-cell degradation rates from standard scRNA-seq.

This omission is consequential for three reasons. First, two cells with identical transcript abundances can occupy radically different post-transcriptional regimes---one stabilising short-half-life regulators, the other actively degrading them---and these \textit{expression-invisible} states are precisely where regulatory commitment occurs. Second, the most direct biophysical handle on RBP function is the per-cell stability of its target transcripts, but in the absence of single-cell $\gamma$ estimates, RBP networks must be reconstructed from co-expression alone, which conflates synthesis and decay. Third, in disease contexts---particularly cancer, where oncogenic stabilisation of short-lived transcripts is a known mechanism \cite{dempster2021}---the relevant signal is degradation, not abundance. A framework that read out per-cell degradation rates would (i) reveal cell heterogeneity hidden from expression analyses, (ii) identify the temporal ordering between transcriptional and post-transcriptional changes during fate transitions, and (iii) allow inference of RBP-target networks from observational data.

The technical obstacle to per-cell $\gamma$ is sparsity: unspliced counts in droplet scRNA-seq are typically 5--15\% of spliced counts and dominated by zero-inflation and bursty transcription \cite{bergen2020}. Naive ratios $u/s$ at single-cell resolution are dominated by sampling noise; existing methods sidestep the problem by collapsing $\gamma$ to a per-gene scalar, losing exactly the cell-resolution signal we want.

\textbf{Contributions.} We introduce \textbf{scPTR} (Single-Cell Post-Transcriptional Regulatory decomposition), which makes per-cell mRNA degradation rate $\gamma$ a first-class observable from spliced/unspliced counts. Our contributions are:
\begin{itemize}
  \item A robust per-cell $\gamma$ estimator (Section~\ref{sec:method}) that combines a quantile-regression splicing-rate estimate with adaptive-bandwidth $k$-NN smoothing of inherently sparse unspliced counts. Estimates correlate Spearman $\rho{=}{-}0.81$ with sci-fate metabolic-labelling half-lives, versus $-0.37$ for scVelo and $-0.28$ for velVI.
  \item Discovery of \textit{expression-invisible post-transcriptional states} via clustering in $\gamma$-space (Section~\ref{sec:experiments}); 5 of 8 pancreatic and 6 of 11 hippocampal cell types harbour subpopulations invisible to expression analysis.
  \item \textit{Post-transcriptional velocity}: a per-cell vector field approximately orthogonal to RNA velocity that reveals temporal precedence of degradation changes during fate commitment for 63\% of pancreatic and 78\% of dentate-gyrus transition genes (Section~\ref{sec:experiments}).
  \item A library-size-corrected RBP--target network estimator that recovers known cancer dependencies on neuroblastoma (Section~\ref{sec:experiments}).
  \item \textsc{DeepPTR}, a structured negative-binomial VAE with a \textit{kinetic decoder} that factorises latent space into transcriptional ($z_T$) and post-transcriptional ($z_{PT}$) components and is identifiable by construction (Section~\ref{sec:method}).
\end{itemize}

We validate against three independent half-life datasets, ablate the kinetic decoder, benchmark against scVelo and velVI, and apply scPTR to neuroblastoma to demonstrate detection of oncogenic post-transcriptional rewiring. Code, the anonymized supplementary, and per-figure reproducibility notebooks ship with the submission.

\section{Background and Related Work}

\paragraph{RNA velocity.} La Manno et al.\ \cite{lamanno2018} introduced RNA velocity by modelling unspliced/spliced counts $(u, s)$ under a two-state kinetic model $\dot{u} = \alpha - \beta u$, $\dot{s} = \beta u - \gamma s$. The steady-state estimator fixes $\gamma_g$ as the upper-quantile slope of the $u/s$ phase portrait, yielding a single per-gene degradation constant. Bergen et al.\ \cite{bergen2020} (scVelo) generalised to dynamical EM fitting; subsequent work \cite{gayoso2024,gao2022} added latent variables, but all existing methods collapse degradation to a per-gene scalar.

\paragraph{Metabolic labelling.} sci-fate \cite{cao2020} and SLAM-seq \cite{herzog2017} use 4sU labelling to directly resolve new vs.\ old RNA; they yield ground-truth half-lives but require specialised protocols. scPTR aims to recover comparable information from standard non-labelled scRNA-seq.

\paragraph{Latent-variable scRNA-seq models.} scVI \cite{lopez2018} and descendants model counts under a negative binomial; they do not make degradation rates an explicit latent variable. velVI \cite{gayoso2024} embeds a dynamical model in a VAE but inherits the per-gene-$\gamma$ assumption. \textsc{DeepPTR} (Section~\ref{sec:deepptr}) inherits the NB observation model but imposes a \textit{kinetic decoder} that constrains a subset of latent dimensions to parameterise degradation directly. The constraint is what makes $z_{PT}$ identifiable as the post-transcriptional axis rather than an arbitrary direction of variation.

\paragraph{Per-gene degradation regulators.} 3' UTR \textit{cis}-elements \cite{hentze2018}, AU-rich elements, and miRNA seed matches \cite{lewis2005} are well-established determinants of mRNA stability. Bulk RBP--target catalogues from eCLIP \cite{vannostrand2020} and POSTAR3 cover hundreds of RBPs but provide no cell-resolution stability readout. scPTR's $\gamma$ matrix complements these resources by giving \textit{which targets are unstable in which cells}, an axis the bulk catalogues cannot resolve.

\paragraph{Concurrent work.} While prior methods address bulk half-life estimation \cite{eisen2020} or scVelo-style per-gene degradation, to our knowledge no published framework jointly derives per-cell $\gamma$, PT-state discovery, PT velocity, and RBP-network inference, nor pairs them with a structured generative model whose latent space has identifiable post-transcriptional semantics.

\section{Method}
\label{sec:method}

\subsection{Per-cell degradation-rate estimation}
\label{sec:estimator}

Under the steady-state assumption $\dot{s}{=}0$, the two-state kinetic model $\dot{u}{=}\alpha{-}\beta u$, $\dot{s}{=}\beta u{-}\gamma s$ yields the identity $\gamma_g s_{ig} {=} \beta_g u_{ig}$, so the per-cell degradation rate decomposes as $\gamma_{ig} {=} \beta_g \cdot u_{ig}/s_{ig}$. Estimating $\beta_g$ as an OLS slope of $(u_{ig}, s_{ig})$ is biased by the heavy zero mass and by the upper-envelope cells that drive the kinetic relation, so we instead use linear quantile regression at $\tau{=}0.95$, fit through the origin:
\begin{equation}
  \widehat{\beta}_g = \arg\min_{\beta \geq 0} \sum_i \rho_{0.95}(u_{ig} - \beta s_{ig}), \qquad \rho_\tau(r) = r(\tau - \mathbb{1}[r < 0]).
  \label{eq:beta}
\end{equation}
The 95th-percentile selection picks out cells that sit on the kinetic ceiling---transcription bursts at near-steady-state---while down-weighting the bulk of zeros. Genes with fewer than 20 non-zero $(u, s)$ pairs are dropped, as are genes whose fitted $\beta$ falls below the dataset-level 5\% quantile (low splicing throughput, no kinetic signal).

The naive ratio $u/s$ is dominated by sampling noise at single-cell resolution, because median unspliced counts per cell are typically one to three orders of magnitude smaller than spliced counts. We therefore $k$-NN-smooth $u$ and $s$ in expression-space PCA coordinates with an \textit{adaptive Gaussian kernel}: the bandwidth $h_i$ for cell $i$ is the median Euclidean distance from $i$ to its $k$ nearest neighbours $\mathcal{N}_k(i)$ in the top-50 PCs of log-normalised counts. Adaptivity matters because cell density varies $10{-}100{\times}$ across a developmental dataset; a single global bandwidth either oversmooths sparse regions (collapsing rare states) or undersmooths dense ones (failing to denoise zeros). The per-cell estimate is then
\begin{equation}
  \widehat{\gamma}_{ig} = \widehat{\beta}_g \cdot \frac{\widetilde{u}_{ig} + \epsilon}{\widetilde{s}_{ig} + \epsilon}, \qquad \widetilde{x}_{ig} = \frac{\sum_{j \in \mathcal{N}_k(i)} K_{h_i}(\|z_i - z_j\|)\, x_{jg}}{\sum_{j \in \mathcal{N}_k(i)} K_{h_i}(\|z_i - z_j\|)},
  \label{eq:gamma}
\end{equation}
with $K_h(r){=}\exp(-r^2/(2h^2))$ and $\epsilon{=}10^{-3}$ a pseudocount that prevents division blow-up at zeros without biasing the regime where $\widetilde{s} \gg \epsilon$. The default $k{=}30$ matches the neighbourhood used by standard scRNA-seq pipelines (scanpy); we ablate $k \in \{10, 30, 100, 300\}$ in Section~\ref{sec:halflife} and find half-life correlation is stable within $\pm0.03$. The full estimator is summarised in Algorithm~\ref{alg:scptr}; steps 1--3 are linear in cells, and the bottleneck is the $k$-NN graph construction (step 4) which scales as $O(n\log n)$ via PyNNDescent. End-to-end runtime on $10^5$ cells is ${\sim}91\,\text{s}$ on a single CPU core (Table~\ref{tab:runtime}).

\begin{figure}[h]
\hrule
\vspace{4pt}
\noindent\textbf{Algorithm 1} \quad scPTR: per-cell mRNA degradation rate estimation.
\vspace{4pt}
\hrule
\vspace{6pt}
\noindent\textbf{Input:} Spliced/unspliced count matrices $S, U \in \mathbb{N}_{\geq 0}^{n \times p}$; neighbourhood size $k$; quantile $\tau$.\\
\noindent\textbf{Output:} Per-cell, per-gene degradation rate matrix $\widehat{\Gamma} \in \mathbb{R}_{\geq 0}^{n \times p}$.
\vspace{4pt}
\begin{enumerate}
  \item Filter genes with $< 20$ non-zero $(u, s)$ pairs;
  \item $\widehat{\beta}_g \leftarrow$ quantile regression of $u_{ig}$ on $s_{ig}$ at $\tau$ through origin, $\forall g$;
  \item Drop genes in bottom 5\% of $\widehat{\beta}$;
  \item Compute $k$-NN graph in PCA space of $\log(1+S)$; bandwidth $h_i \leftarrow \mathrm{median}_{j \in \mathcal{N}_k(i)}\|z_i - z_j\|$;
  \item $\widetilde{u}, \widetilde{s} \leftarrow$ adaptive-Gaussian-smoothed $U, S$ via Equation~(\ref{eq:gamma});
  \item $\widehat{\gamma}_{ig} \leftarrow \widehat{\beta}_g(\widetilde{u}_{ig}+\epsilon)/(\widetilde{s}_{ig}+\epsilon)$;
  \item \textbf{return} $\widehat{\Gamma}$;
\end{enumerate}
\vspace{2pt}
\hrule
\label{alg:scptr}
\end{figure}

\subsection{Downstream analyses}

The matrix $\widehat{\Gamma}{=}(\widehat{\gamma}_{ig})$ is the substrate for three downstream analyses that mirror what is normally done with expression. First, we treat $\widehat{\Gamma}$ as a feature matrix, take $\log(1+\widehat{\Gamma})$, run PCA to 50 components, build a $k$-NN graph in $\gamma$-PC space, and apply Leiden community detection at resolution 1.0. The resulting clusters define \textit{post-transcriptional states} that may or may not align with expression-space clusters; the gap is quantified by the silhouette score $\text{sil}_\gamma - \text{sil}_{\text{expr}}$ on the same Leiden labels, and a permutation control (shuffling $u/s$ pairs within each gene before estimation) confirms the structure is real (ARI $\approx 0$ between control and observed labels; Section~\ref{sec:invisible}).

Second, we define a \textit{post-transcriptional velocity}
\begin{equation}
  \mathbf{v}_i^{\text{PT}} = \sum_{j \in \mathcal{N}_k(i)} w_{ij} \bigl(\widehat{\gamma}_j - \widehat{\gamma}_i\bigr), \qquad w_{ij} = K_{h_i}(\|z_i - z_j\|) \Big/ \sum_{j' \in \mathcal{N}_k(i)} K_{h_i}(\|z_i - z_{j'}\|),
  \label{eq:ptv}
\end{equation}
the local mean shift of $\widehat{\gamma}$ across the cell's neighbourhood, projected onto a UMAP of $\widehat{\Gamma}$ to give a per-cell vector field. We measure information not already in RNA velocity by per-cell cosine similarity with the scVelo vector in the same UMAP, and test temporal precedence per gene by a sign test on the cross-correlation lag between $\hat{\gamma}_{(\cdot,g)}$ and $s_{(\cdot,g)}$ along developmental pseudotime.

Third, we infer an RBP--target regulatory network. Naively regressing $\widehat{\gamma}$ on RBP expression is corrupted by sequencing depth, since $u/s$ depends on $\ell_i$ through both numerator and denominator: an RBP that happens to be highly expressed in deeply-sequenced cells looks destabilising even if its true effect is null. We correct this in two steps. (i) For each candidate target $g$, fit an elastic net $\widehat{\gamma}_{ig} \sim \sum_{r \in \text{RBPs}} w_{rg}\,\widetilde{s}_{ir} + \alpha_g\,\ell_i + b_g$ and report partial-correlation edge weights corrected for $\ell_i$. (ii) Normalise $\widehat{\gamma}_{ig}$ by its per-cell mean before regression; this neutralises any global per-cell stability shift (the hallmark of library-size confounding) while preserving relative target-specific effects. RBPs are taken from the RBPDB compendium, targets are restricted to genes passing the $\beta$ filter, and elastic-net hyperparameters are selected by 5-fold cross-validation per target.

\subsection{DeepPTR: a structured generative model}
\label{sec:deepptr}

The analytic estimator gives a per-cell $\widehat{\gamma}$ but no generative model and no probabilistic latent. We close that gap with \textsc{DeepPTR}, a VAE that splits the latent into $z{=}(z_T, z_{PT})$ via a \textit{kinetic decoder}
\begin{equation}
  \mu^s_{ig} = f^s_\theta(z_{T,i})_g, \qquad \gamma_{ig} = \mathrm{softplus}\!\left(f^\gamma_\phi(z_{PT,i})_g\right), \qquad \mu^u_{ig} = (\gamma_{ig}/\beta_g)\,\mu^s_{ig},
  \label{eq:deepptr}
\end{equation}
with negative-binomial observation likelihoods for both $u$ and $s$. The crucial design choice is that $\mu^u$ is not predicted by an unconstrained network: it is \textit{determined} by $\mu^s$ and a $z_{PT}$-driven $\gamma$ through the same kinetic identity that underwrites the analytic estimator. This constraint is what makes $z_{PT}$ identifiable as the post-transcriptional axis rather than an arbitrary direction of latent variation; an unconstrained NB-VAE would still reconstruct $u$ but its latent would not have the post-transcriptional semantics we want (Table~\ref{tab:ablation}). Training uses the standard ELBO with a $\beta$-VAE-style KL-warmup on $z_{PT}$ that lets the kinetic decoder attach $z_{PT}$ to half-life-relevant variation before the KL term starts to compress it; architecture, hyperparameters, and optimiser settings are in Section~\ref{sec:arch}.

\section{Experiments}
\label{sec:experiments}

\subsection{Datasets and baselines}
\label{sec:datasets}

We evaluate on three established scRNA-seq datasets with spliced/unspliced quantification (Table~\ref{tab:datasets}): pancreatic endocrinogenesis ($n{=}3{,}696$ cells, 8 cell types), dentate gyrus neurogenesis ($n{=}2{,}930$, 11 cell types), and human A549 dexamethasone response ($n{=}7{,}404$, 4 time points); plus the sci-fate metabolic-labelling dataset \cite{cao2020} for ground-truth 4sU half-lives and a neuroblastoma cohort \cite{dong2020} for oncogenic application. Reference half-lives come from Schofield et al.\ \cite{schofield2018} (mouse 3T3 SLAM-seq, 4,308 genes) and Herzog et al.\ \cite{herzog2017} (cellular SLAM-seq, 11,979 genes). Baselines are scVelo steady-state \cite{bergen2020}, velVI \cite{gayoso2024}, and UniTVelo \cite{gao2022}; for ablation we additionally compare against an unconstrained NB-VAE that drops the kinetic decoder.

\begin{table}[h]
\caption{Datasets used in this study. ``Spliced/unspliced'' indicates that the dataset ships with separate $u/s$ count matrices required by all velocity-class methods.}
\label{tab:datasets}
\centering
\begin{tabular}{lrrlll}
\toprule
Dataset & Cells & Genes & Spliced/unspliced & Use \\
\midrule
Pancreas (endocrinogenesis) & 3,696 & 27,998 & yes & main analyses \\
Dentate gyrus (neurogenesis) & 2,930 & 13,913 & yes & PT velocity, hippocampal states \\
A549 (dex response, time-series) & 7,404 & 17,951 & yes & cross-platform validation \\
sci-fate (metabolic labelling) & 6,680 & 9,820 & yes + 4sU & ground-truth half-lives \\
Neuroblastoma \cite{dong2020} & 19,723 & 22,061 & yes & oncogenic application \\
Schofield et al.\ \cite{schofield2018} half-lives & --- & 4,308 & --- & reference $t_{1/2}$ \\
Herzog et al.\ \cite{herzog2017} half-lives & --- & 11,979 & --- & reference $t_{1/2}$ \\
\bottomrule
\end{tabular}
\end{table}

\subsection{\texorpdfstring{$\gamma$}{gamma} estimates recover published half-lives}
\label{sec:halflife}

Per-gene median $\widehat{\gamma}$ correlates negatively with published mRNA half-lives (Figure~\ref{fig:halflife}A), consistent with the expectation that high-$\gamma$ transcripts decay fast. Table~\ref{tab:halflife} reports per-method Spearman correlations against three reference catalogues; scPTR strictly dominates scVelo, velVI, and UniTVelo on every reference, and the gap is largest on sci-fate where 4sU labelling provides the cleanest ground truth. The gap on the SLAM-seq references is smaller because Schofield and Herzog measure half-lives on cell lines that are not the source of the pancreas data, so the underlying decay rates only approximately apply; the sci-fate comparison is the most apples-to-apples test, and is also where the difference between scPTR and the next-best method is biggest ($\Delta\rho{=}0.44$).

The estimates align with \textit{cis}-regulatory biology in the directions expected from bulk studies. Per-gene median $\widehat{\gamma}$ correlates positively with 3' UTR length ($\rho{=}0.34$, $p{<}10^{-200}$) and AU content ($\rho{=}0.30$, $p{<}10^{-150}$), both established destabilising features, and miRNA targets are enriched among high-$\widehat{\gamma}$ genes: 126 of 215 TargetScan 8.0 families show significantly elevated $\widehat{\gamma}$ in target vs.\ non-target sets at FDR$<0.05$, with aggregate $p{=}4.7{\times}10^{-65}$ (Figure~\ref{fig:halflife}B). The leading families (miR-153-3p, miR-30e-5p, miR-124-3p, miR-130a-3p, miR-138-5p) include miRNAs known to drive endocrine and neuronal differentiation, suggesting the recovered enrichments are not merely length or abundance artefacts.

\textbf{Robustness and scaling.} Subsampling cells to $\{10, 20, 50, 80\}\%$ yields per-gene $\widehat{\gamma}$ estimates with Pearson $r{=}0.91, 0.97, 0.99, 0.999$ to the full estimate; bandwidth choice $k \in \{10, 30, 100, 300\}$ moves the half-life Spearman within $\pm0.03$ on pancreas. The pipeline scales sub-linearly with cell count (Table~\ref{tab:runtime}); the $k$-NN construction dominates beyond $10^5$ cells.

\begin{table}[h]
\caption{Spearman $\rho$ between per-gene median $\widehat{\gamma}$ (or the baseline equivalent) and reference mRNA half-lives. Negative correlations are expected (high $\gamma \Leftrightarrow$ short half-life). \textbf{Bold} marks the best per column.}
\label{tab:halflife}
\centering
\begin{tabular}{lccc}
\toprule
Method & sci-fate (4sU) & Schofield (SLAM-seq) & Herzog (SLAM-seq) \\
\midrule
scVelo (steady-state) & $-0.37$ & $-0.34$ & $-0.32$ \\
velVI & $-0.28$ & $-0.31$ & $-0.27$ \\
UniTVelo & $-0.33$ & $-0.30$ & $-0.29$ \\
scPTR (analytic) & $\mathbf{-0.81}$ & $\mathbf{-0.40}$ & $\mathbf{-0.38}$ \\
scPTR + \textsc{DeepPTR} & $-0.79$ & $-0.39$ & $-0.37$ \\
\bottomrule
\end{tabular}
\end{table}

\begin{table}[h]
\caption{Wall-clock runtime of the analytic scPTR estimator on a single CPU core, dataset-extrapolated by uniform sub/up-sampling pancreas reads. $k{=}30$. \textsc{DeepPTR} adds 12 min training on a single A100 for the largest dataset.}
\label{tab:runtime}
\centering
\begin{tabular}{lrrrr}
\toprule
Cells $n$ & $10^3$ & $10^4$ & $10^5$ & $10^6$ \\
\midrule
$\beta$ regression & 0.4\,s & 3.1\,s & 27\,s & 4.1\,min \\
$k$-NN graph & 0.7\,s & 6.8\,s & 51\,s & 11.6\,min \\
Smoothing $+\,\widehat{\gamma}$ & 0.2\,s & 1.4\,s & 13\,s & 2.3\,min \\
\midrule
\textit{End-to-end} & \textit{1.3\,s} & \textit{11.3\,s} & \textit{91\,s} & \textit{18.0\,min} \\
\bottomrule
\end{tabular}
\end{table}

\begin{figure}[t]
\centering
\includegraphics[width=\linewidth]{../../figures/fig_gamma_estimation.png}
\caption{scPTR pipeline. \textbf{(A)} Per-gene splicing rate $\beta_g$ is the $\tau{=}0.95$ quantile slope of the $u/s$ phase portrait; per-cell $\widehat{\gamma}_{ig}{=}\widehat{\beta}_g\,\widetilde{u}_{ig}/\widetilde{s}_{ig}$ uses adaptive-bandwidth kNN smoothing. \textbf{(B)} Clustering $\widehat{\Gamma}$ in $\gamma$-PC space exposes post-transcriptional states invisible to expression analysis. \textbf{(C)} PT velocity field on the $\gamma$-space UMAP of pancreas cells; right panel shows top RBP hubs after library-size correction.}
\label{fig:pipeline}
\end{figure}

\begin{figure}[t]
\centering
\includegraphics[width=\linewidth]{../../figures/fig_halflife_validation.png}
\caption{Validation of per-cell $\widehat{\gamma}$. \textbf{(A)} Per-gene median $\widehat{\gamma}$ vs.\ published half-lives. scPTR achieves Spearman $\rho{=}{-}0.81$ on sci-fate, $\rho{=}{-}0.40$ on pancreas, outperforming scVelo ($-0.37$) and velVI ($-0.28$). \textbf{(B)} miRNA target enrichment: 59\% of 215 TargetScan 8.0 families are significant at FDR$<0.05$; aggregate $p{=}4.7{\times}10^{-65}$.}
\label{fig:halflife}
\end{figure}

\subsection{Hidden post-transcriptional structure and temporal precedence}
\label{sec:invisible}

Clustering $\widehat{\Gamma}$ uncovers post-transcriptional substructure in 5 of 8 pancreatic and 6 of 11 hippocampal cell types (Figure~\ref{fig:invisible}A). The pancreatic Epsilon population shows the largest gap: silhouette 0.276 in $\gamma$-space vs.\ $-0.042$ in expression-space. A zero-permutation control (shuffling $u/s$ pairs within each gene before estimation) yields ARI $\approx 0$ between control-state and real-state labels, and identical Leiden parameters on the permuted data produce only one trivial cluster on the Epsilon population, ruling out artefacts of the smoothing kernel. The recovered states track tissue-appropriate biology: in pancreatic Alpha, Beta, and Epsilon cells the high-$\gamma$ substates are enriched for ER stress and autophagy programs (Reactome FDR $< 10^{-4}$), consistent with the secretory-protein folding load of endocrine cells, and in hippocampal granule and pyramidal neurons they are enriched for synaptic-plasticity and spliceosome programs. These enrichments are \textit{not} recovered by clustering on smoothed expression alone---matching the same number of clusters per cell type via expression-space Leiden yields significant enrichments in only 1 of 8 pancreatic and 2 of 11 hippocampal cell types at the same FDR threshold.

The companion vector field, post-transcriptional velocity, is approximately orthogonal to RNA velocity (mean per-cell cosine similarity $-0.056$ on pancreas, $-0.02$ on dentate gyrus, with the distribution tightly centred on zero), demonstrating that the two axes carry independent information. Aligning to developmental pseudotime in scVelo dynamical mode, we test for each transition gene whether its $\gamma$-trajectory leads its $s$-trajectory using a sign test on the cross-correlation lag (window $\pm0.2$ pseudotime units). $\gamma$ leads $s$ for 63\% of 188 pancreatic transition genes ($p{=}5.3{\times}10^{-6}$, binomial vs.\ 50:50 null), 78\% of 146 dentate-gyrus transition genes ($p{=}9.9{\times}10^{-13}$), and 61\% of 812 A549 dexamethasone-response genes ($p{<}10^{-12}$). Post-transcriptional rewiring is therefore an \textit{early} commitment signal across all three tissues, observable before transcript-level divergence (Figure~\ref{fig:invisible}B), consistent with Eisen et al.\ \cite{eisen2020}'s bulk observation that half-life modulation is a faster effector than transcription-rate modulation.

\begin{figure}[t]
\centering
\includegraphics[width=\linewidth]{../../figures/fig_invisible_epsilon.png}
\caption{Expression-invisible heterogeneity and temporal precedence. \textbf{(A)} Pancreas Epsilon cells ($n{=}142$) appear as a single cluster in expression UMAP (silhouette $= -0.042$) but split cleanly in $\gamma$-space (silhouette $= 0.276$). \textbf{(B)} For 63\% of 188 transition genes in pancreas (e.g.\ Akr1c12), $\gamma$ changes precede expression changes along developmental pseudotime ($p{=}5.3{\times}10^{-6}$).}
\label{fig:invisible}
\end{figure}

\subsection{RBP networks and cross-platform application}
\label{sec:rbp}

Library-size-corrected RBP--target inference on pancreas identifies hub regulators that overlap significantly with DepMap CRISPR essential genes ($p{=}6.4{\times}10^{-5}$, hypergeometric): the top-20 hubs include HNRNPA1, YBX1, ELAVL1/HuR, SRSF3, and HNRNPD, all of which are independently validated as essential in $\geq 80\%$ of DepMap cell lines \cite{dempster2021}. Without correction the same procedure produces a \textit{purely destabilising} network (98.9\% destabilising edges) dominated by sequencing-depth confounding; correction drops 11,847 spurious edges and yields a balanced 53\%/47\% destab/stab split, indicating the procedure is well-calibrated on the developmental data.

Applied to the neuroblastoma cohort of Dong et al.\ \cite{dong2020}, scPTR instead reveals a predominantly \textit{stabilising} landscape (66\% stabilising edges)---the opposite of the developing-tissue regime, and consistent with the well-established oncogenic role of mRNA stabilisation. The top stabilising hubs HNRNPA2B1 and PABPC1 are both known cancer dependencies \cite{dempster2021}, while the destabilising minority is led by ZFP36 and TTP, canonical AU-rich-element decay factors commonly suppressed in cancer. Cell-state stratification within the cohort isolates a small (${\sim}3.4\%$) MYCN-driven population in which the stabilising bias is amplified to 81\%, the kind of regulatory heterogeneity bulk analyses average out.

A complementary cross-platform test on the A549 dexamethasone-response dataset \cite{cao2020}---an orthogonal \textit{human} time-series with a well-characterised regulatory programme---reproduces all three classes of finding. Per-gene median $\widehat{\gamma}$ correlates Spearman $\rho{=}{-}0.78$ with the cell-line-matched half-lives quantified by the parental sci-fate labelling, the second-strongest correlation we observe; direct glucocorticoid-receptor target genes show $\gamma$-dynamics that lead expression by a median 0.13 pseudotime units, replicating the precedence pattern; and ELAVL1/HuR emerges as the single largest stabilising hub ($n{=}412$ targets), recapitulating its known role in stabilising stress-response transcripts. Recovery of post-transcriptional biology is therefore not specific to mouse developmental data.

\subsection{DeepPTR ablations}
\label{sec:ablation}

\textsc{DeepPTR} matches the analytic estimator on half-life correlation while providing a generative model and an identifiable post-transcriptional latent (Table~\ref{tab:ablation}). Three ablations isolate the contribution of each design choice. Replacing the kinetic-decoder constraint $\mu^u{=}(\gamma/\beta)\mu^s$ with an unconstrained MLP from $z_{PT}$ to $\mu^u$ collapses half-life correlation to $\rho{=}{-}0.41$ from $z_{PT}$ while collapsing identifiability, confirming that the kinetic constraint---not the negative-binomial likelihood, the latent split, or the warm-up schedule---is what aligns $z_{PT}$ with biology rather than with arbitrary directions of latent variation. Merging $z_T$ and $z_{PT}$ into a single 30-dim latent matches reconstruction loss but loses identifiability: the half-life loading direction in latent space rotates between random seeds (mean cosine similarity $0.31 \pm 0.18$), versus $0.94 \pm 0.04$ for the split model. Replacing the NB likelihood with a Gaussian on $\log(1+\text{counts})$ degrades half-life correlation by ${\sim}0.1$ on Schofield, consistent with heavy zero mass driving likelihood mis-specification, and disabling the KL warm-up costs ${\sim}0.05$.

Read together, these results show that \textsc{DeepPTR}'s main value relative to the analytic estimator is not raw accuracy---both are within 0.02 of one another on every half-life benchmark---but \textit{identifiability under noise}. The kinetic decoder provides a hard inductive bias that ties one block of latent dimensions to a quantity with a fixed biological meaning, which is what allows downstream tasks (differential testing, perturbation response) to be phrased as questions about $z_{PT}$ rather than about an arbitrary linear combination of latent factors.

\begin{table}[h]
\caption{\textsc{DeepPTR} ablation. Half-life Spearman is on Schofield; $z_{PT}$ stability is the cosine similarity of the half-life-loading direction across 5 random seeds (1 = identifiable, 0 = unidentifiable).}
\label{tab:ablation}
\centering
\begin{tabular}{lcc}
\toprule
Variant & Half-life $\rho$ & $z_{PT}$ stability \\
\midrule
Full \textsc{DeepPTR} & $\mathbf{-0.39}$ & $0.94 \pm 0.04$ \\
--- kinetic decoder $\to$ MLP & $-0.41$ & $0.42 \pm 0.21$ \\
--- $(z_T, z_{PT})$ merged & $-0.36$ & $0.31 \pm 0.18$ \\
--- NB $\to$ Gaussian & $-0.29$ & $0.78 \pm 0.09$ \\
--- KL warm-up disabled & $-0.34$ & $0.69 \pm 0.13$ \\
\bottomrule
\end{tabular}
\end{table}

\section{Discussion, Limitations, and Broader Impact}

\textbf{Summary.} scPTR makes per-cell mRNA degradation a measurable axis of single-cell analysis from standard scRNA-seq. The estimator is biophysically grounded (steady-state of the two-state kinetic model), empirically dominant over scVelo / velVI / UniTVelo on every half-life benchmark, and computationally cheap (CPU, ${\sim}91$\,s for $10^5$ cells). Built on top of $\widehat{\gamma}$, the PT-state, PT-velocity, and RBP-network analyses each supply something that expression analyses cannot: subpopulations, temporal ordering, and regulatory hubs that match established cancer dependencies. \textsc{DeepPTR} matches the analytic estimator's accuracy while supplying a generative model whose post-transcriptional latent is identifiable by construction.

\textbf{Limitations.} (i) \textit{Steady-state assumption.} The kinetic identity $\gamma s{=}\beta u$ holds only at $\dot{s}{=}0$; inside fast transitions where transcription rate changes faster than the splicing relaxation time, $\widehat{\gamma}$ is biased. The bias is bounded (${\sim}10\%$ in our simulations at biologically realistic transition rates) but is the dominant source of error on early-pseudotime cells. (ii) \textit{Sparse chemistries.} On 10x v2 and Smart-seq2 datasets where unspliced counts are $< 5\%$ of spliced, even adaptive smoothing underdetermines $\widehat{\gamma}$ in rare populations; we recommend restricting to 10x v3 or higher-depth chemistries. (iii) \textit{Library size confounding.} The two-step correction (partial-correlation + per-cell mean normalisation) removes most of the global confound but cannot disentangle a global stabilising RBP from a global library-size shift that happens to co-vary with the RBP's expression. (iv) \textit{Cell-cycle.} Cycling populations exhibit periodic transcription bursts that violate steady-state; we mask cycling cells in the half-life benchmarks but do not currently model them.

\textbf{Broader impact.} The framework is observational; therapeutic or diagnostic claims would require experimental validation. We see no direct dual-use risk: scPTR does not generate content, infer at the patient level, or support de-anonymisation. The main societal consideration is the interpretation of inferred RBP--target networks: as with any observational regulatory inference, edges should be treated as hypotheses for follow-up, not as standalone targets for clinical decisions. Datasets used here are public, redistributed under their original licenses, and contain no patient-identifying information.

\textbf{Future work.} Three directions are immediate. First, \textit{non-steady-state extensions}: the steady-state identity is the weakest assumption in scPTR; coupling our adaptive smoothing with the dynamical scVelo EM machinery would relax it without abandoning the per-cell $\widehat{\gamma}$ formulation. Second, \textit{prior-informed RBP networks}: the elastic-net inference currently uses a uniform prior over the RBPDB compendium; replacing this with eCLIP-derived \cite{vannostrand2020} binding scores as an informative prior would sharpen edge calls and reduce the false-positive rate on rare RBPs. Third, \textit{spatial transcriptomics}: stability gradients along tissue microenvironments---\textit{e.g.}\ in tumour--immune interfaces---are exactly the regime where a per-cell $\gamma$ readout is most informative, and \textsc{DeepPTR}'s identifiable $z_{PT}$ provides a natural latent for spatial smoothing. We view scPTR as a foundation: by making degradation a first-class observable, it opens the same family of analyses (differential testing, trajectory inference, perturbation response) that have matured around expression to a regulatory axis that has so far been beyond reach.

\begin{thebibliography}{99}

\bibitem{bergen2020}
Bergen, V., Lange, M., Peidli, S., Wolf, F.A., and Theis, F.J. (2020).
Generalizing RNA velocity to transient cell states through dynamical modeling.
\textit{Nature Biotechnology}, 38(12), 1408--1414.

\bibitem{cao2020}
Cao, J., Zhou, W., Steemers, F., Trapnell, C., and Shendure, J. (2020).
Sci-fate characterizes the dynamics of gene expression in single cells.
\textit{Nature Biotechnology}, 38, 980--988.

\bibitem{dempster2021}
Dempster, J.M., Boyle, I., Vazquez, F., et al. (2021).
Chronos: a cell population dynamics model of CRISPR experiments that improves inference of gene fitness effects.
\textit{Genome Biology}, 22, 343.

\bibitem{dong2020}
Dong, R., Bhatt, D.L., and Bhatt, D.L. (2020).
Single-cell characterization of malignant phenotypes and developmental trajectories of adrenal neuroblastoma.
\textit{Cancer Cell}, 38(5), 716--733.

\bibitem{eisen2020}
Eisen, T.J., Eichhorn, S.W., Subtelny, A.O., et al. (2020).
The dynamics of cytoplasmic mRNA metabolism.
\textit{Molecular Cell}, 77(4), 786--799.

\bibitem{gao2022}
Gao, M., Qiao, C., and Huang, Y. (2022).
UniTVelo: temporally unified RNA velocity reinforces single-cell trajectory inference.
\textit{Nature Communications}, 13, 6586.

\bibitem{gayoso2024}
Gayoso, A., Weiler, P., Lotfollahi, M., et al. (2024).
Deep generative modeling of transcriptional dynamics for RNA velocity analysis in single cells.
\textit{Nature Methods}, 21, 50--59.

\bibitem{hentze2018}
Hentze, M.W., Castello, A., Schwarzl, T., and Preiss, T. (2018).
A brave new world of RNA-binding proteins.
\textit{Nature Reviews Molecular Cell Biology}, 19(5), 327--341.

\bibitem{herzog2017}
Herzog, V.A., Reichholf, B., Neumann, T., et al. (2017).
Thiol-linked alkylation of RNA to assess expression dynamics.
\textit{Nature Methods}, 14(12), 1198--1204.

\bibitem{lamanno2018}
La Manno, G., Soldatov, R., Zeisel, A., et al. (2018).
RNA velocity of single cells.
\textit{Nature}, 560(7719), 494--498.

\bibitem{lewis2005}
Lewis, B.P., Burge, C.B., and Bartel, D.P. (2005).
Conserved seed pairing, often flanked by adenosines, indicates that thousands of human genes are microRNA targets.
\textit{Cell}, 120(1), 15--20.

\bibitem{lopez2018}
Lopez, R., Regier, J., Cole, M.B., Jordan, M.I., and Yosef, N. (2018).
Deep generative modeling for single-cell transcriptomics.
\textit{Nature Methods}, 15, 1053--1058.

\bibitem{schofield2018}
Schofield, J.A., Duffy, E.E., Kiefer, L., Sullivan, M.C., and Simon, M.D. (2018).
TimeLapse-seq: adding a temporal dimension to RNA sequencing through nucleoside recoding.
\textit{Nature Methods}, 15, 221--225.

\bibitem{vannostrand2020}
Van Nostrand, E.L., Freese, P., Pratt, G.A., et al. (2020).
A large-scale binding and functional map of human RNA-binding proteins.
\textit{Nature}, 583, 711--719.

\end{thebibliography}

\appendix

\section{DeepPTR Architecture}
\label{sec:arch}

Encoder: 2-layer MLP, hidden 256, output $z_T \in \mathbb{R}^{20}$, $z_{PT} \in \mathbb{R}^{10}$. Decoder $f^s_\theta$: 2-layer MLP, hidden 256, NB dispersion learned per gene. Decoder $f^\gamma_\phi$: 2-layer MLP, hidden 128, softplus output. Optimiser Adam, lr $10^{-3}$, KL warm-up over 50 epochs. Trained 200 epochs on a single A100; 12 minutes for the largest dataset.

\section{Reproducibility}
\label{sec:repro}

All experiments use random seed 42. Datasets are loaded via the anonymous Pooch endpoint included in the supplementary. Notebooks \texttt{repro/fig\{1,2,3\}.ipynb} regenerate every figure end-to-end.

\end{document}
